%% Introduction %%

Nous allons dès à présent appuyer sur le métier de documentaliste et sa relation avec les collections et bases de données ainsi que sur l'importance de leur expertise dans la médiation des collections.

%% Situer dans le temps %%

Il va de soi que les premières personnes concernées par les bases de données, les données de descriptions et le \gls{td} appliqués sont les personnes travaillant au sein des institutions patrimoniales. C'est en partie à travers leurs besoins, leur expertise et leurs recommandations que les politiques documentaires sont élaborées. Pour autant, l’accessibilité aux collections reste un enjeu de tous les jours. 

\subsection{Connaissances de outils de consultation et de production des bases de données}

%% Habitudes, données de description préférées %%

Au fur et à mesure des années, des changements dans les bases de données et de la mise en place de nouveaux outils, les professionnels de l'INA ont construits leurs habitudes sur ces changements. Quand les archives professionnelles et le dépôt légal étaient séparés, les outils de consultation des bases de données l'étaient aussi. Une fois les services réunis, les documentalistes conservent majoritairement leurs habitudes. La documentaliste que nous avons interrogé, anciennement au dépôt légal, continue majoritairement à utiliser \gls{hy}, l'outil de consultation du dépôt légal mais dans le cas de reprise d'antériorité, elle doit utiliser \gls{to}, l'outil de consultation et de production des archives professionnelles, duquel dépendent les reprises d'antériorité. Ces professionnels sélectionnent donc leurs outils en fonction de leurs habitudes mais également en fonction des missions. Malgré la réunification des services des archives professionnelles et du dépôt légal, les deux outils ont des usages différents. La transition vers un nouvel outil peut être difficile, comme l'est actuellement la transition vers l'outil de consultation Notilus\renvoientretien{\DA}{les bases de données de l'INA}. Ce dernier n'est pour le moment qu'à usage interne, il a été mis en place dans des contextes différents d'\gls{hy} et \gls{to}. Il contient toutes les transcriptions réalisées par IA dans lesquelles il est possible de faire des recherches tout-texte ainsi qu'une recherche sémantique, ces deux fonctionnalités pouvant compenser le manque d'informations des données de description comme c'est le cas dans les fonds de l'ORTF\renvoientretien{\PL}{les bases de données de l'INA}.

Une part du métier reste l'adaptabilité. Ces experts doivent pouvoir répondre à toutes les questions des chercheurs, quel que soit leur sujet. De plus, ayant connaissance des collections de l'INA, ils doivent pouvoir faire appel à chaque outils à leur disposition. C'est pourquoi ils ont l'habitude d'utiliser ces trois outils de consultations mais certains font également appel aux fiches papier ou aux archives physiques. Quand une information dans la base de données est incohérente, la réponse peut se trouver dans un programme TV par exemple\renvoientretien{\PL}{l'expertise du documentaliste}. 

\subsection{L'expertise du documentaliste}

Si les chercheurs ont un sujet privilégié qu'ils vont explorer avec l'aide des documentalistes. Ces dernièr.es doivent avoir une connaissance parfaite des fonds, c'est \enquote{l'expertise du documentaliste}. Ces professionnels ont une connaissance poussée sur le fonctionnement des bases de données, sur les différentes politiques documentaires, sur la construction et l'acquisition des collections. Cela ne concerne évidemment pas que l'INA, les documentalistes acquièrent une expertise dans le domaine dans lequel ils exercent. C'est pourquoi, dans le cadre de commandes, les documentalistes de l'INA peuvent faire appel à des documentalistes en dehors de l'INA, dans le cas où les documents nécessaires ne sont pas dans leurs fonds mais dans les fonds d'une autre institution.

Pourtant, c'est aujourd'hui une connaissance qui est en péril. Il y a plus de 30 ans, le dépôt légal des documents audiovisuels chamboulait le fonctionnement de l'INA, à cette occasion, plusieurs documentalistes ont été engagé.es pour accompagner ce nouveau flux entrant de documents. Ce qui signifie qu'actuellement, les documentalistes ayant vu naître et grandir ce dépôt légal approchent de la retraite ou y sont déjà. Le \gls{td} n'ayant pas ou peu été documenté, il est plus que jamais essentiel qu'ils.elles transmettent leurs connaissances comme leurs prédécesseuses l'ont fait. C'est une démarche déjà entamée pour leur connaissance des fonds. Dans le cas où ces professionnelles investissent beaucoup de temps sur un fond, sur une typologie ou sur un sujet, elles en deviennent des expertes et sont en capacité de transmettre leurs connaissances à travers les corpus qu'elles construisent. Nous retrouvons ces corpus dans l'outil de consultation et de production \enquote{\gls{to}} ou sur les différentes plateformes de médiation de l'INA\renvoientretien{\PL}{l'expertise du documentaliste}. 

Malheureusement, leurs connaissances sur le \gls{td} ne sont pas aussi bien documentées, ce qui est un point important à mettre en avant dans la construction de notre outil de cartographie du \gls{td}. De cette façon, l'accessibilité aux bases de données de l'INA par les documentalistes sont soumises aux mêmes enjeux que les chercheurs à la différence qu'ils possèdent toutes les clefs pour comprendre et utiliser les bases de données de l'institut. 